\subsection*{Impactos tecnológicos}
\addcontentsline{toc}{subsection}{Impactos tecnológicos}

El desarrollo de la plataforma representa un avance significativo en la aplicación de arquitecturas de software modernas para la resolución de problemas socioeconómicos. Desde la perspectiva de la Ingeniería de Sistemas, el impacto tecnológico se manifiesta en la implementación de una infraestructura escalable, segura y adaptativa, cuyo alcance se define en tres horizontes temporales.

En el corto plazo, el impacto principal radica en la validación de una arquitectura basada en microservicios y la implementación de contenedores mediante Docker. Esta aproximación garantiza una alta disponibilidad del servicio y permite ciclos de despliegue continuo (CI/CD) ágiles, asegurando que las herramientas de simulación y las calculadoras financieras operen con tiempos de respuesta óptimos para el usuario final. Un componente fundamental en esta etapa es la flexibilidad metodológica; al tratarse de una startup, la arquitectura técnica está diseñada para permitir cambios rápidos y adaptaciones basadas en las sugerencias de los usuarios o ajustes en el modelo de negocio. Esta capacidad de iteración técnica asegura que el producto evolucione en sintonía con las necesidades reales del mercado sin comprometer la estabilidad del sistema.

A mediano plazo, el impacto tecnológico se expande hacia la consolidación de modelos de análisis de datos y la personalización avanzada de la experiencia educativa. Mediante la implementación de algoritmos de machine learning, la plataforma logra generar rutas de aprendizaje adaptativas que se ajustan al perfil de riesgo, nivel de conocimiento y comportamiento financiero de cada usuario. El uso de infraestructuras cloud bajo demanda permite que estos modelos de procesamiento crezcan horizontalmente conforme aumenta la base de usuarios en Bogotá, sin degradar el rendimiento del aplicativo. En esta fase, la plataforma se posiciona como un nodo de innovación al integrar protocolos de seguridad de la información robustos, garantizando la integridad y confidencialidad de los datos de acuerdo con los estándares de la Ley 1581 de 2012 y las exigencias de ciberseguridad del sector Fintech.

Finalmente, a largo plazo, se proyecta un impacto transformador mediante la maduración de un ecosistema de Open Finance. La arquitectura modular de la plataforma facilitará la interoperabilidad con otras entidades del sector mediante APIs estandarizadas, promoviendo un entorno de colaboración tecnológica que reduce las barreras de entrada para nuevos servicios financieros. El éxito de este modelo tecnológico servirá como referente para el desarrollo de soluciones similares en otras regiones, demostrando que la ingeniería de software es el motor fundamental para la sostenibilidad operativa y la escalabilidad de los modelos de negocio digitales en Colombia.